\documentclass{article}
\usepackage[utf8]{inputenc}
\usepackage{geometry}
\usepackage{fancyhdr}
\usepackage{enumitem}
\usepackage{amsmath}
\usepackage{amssymb}
\usepackage{xcolor}
\usepackage{multicol}
\usepackage{xfrac}
\usepackage{parskip}
\usepackage{tikz}
\setcounter{MaxMatrixCols}{20}

\geometry{margin=1in}
\pagestyle{fancy}
\fancyhf{}
\rhead{210 Notes}
\lhead{Class: Math 210, Sacramento State}
\cfoot{\thepage}

\title{Abstract Algebra Comprehensive Exam Practice}
\author{Ryan Harris}
\date{\today}

\begin{document}

\maketitle

\setcounter{section}{-1}
\section{Introduction}

\newpage

\section{The Sylow Theorems}

\begin{enumerate}
    \item Let $G$ be a finite abelian group and suppose $p$ divides $|G|$. Prove that $n_p = 1$.\\
    \textbf{Proof:} Let $G$ be a finite abelian group and suppose that $p \big| |G|$. Then we have
    \begin{align*}
        p \big| |G| &\Longrightarrow \exists m \in \mathbb{Z}^+:pm=|G| \\
        &\Longrightarrow \exists P \in Syl_p(G) \\
        &\Longrightarrow n_p \geq 1
    \end{align*}
    Now, assume $\exists Q \in Syl_p(G)$ such that $P \neq Q$. Then we have
    \begin{align*}
        \exists Q \in Syl_p(G): P \neq Q &\Longrightarrow \exists g \in G: Q \leq gPg^{-1} \\
        &\Longrightarrow q = gpg^{-1} \text{ with }q \in Q, p \in P \\
        &\Longrightarrow qg = gp \\
        &\Longrightarrow gq = gp \text{ since } G \text{ is abelian} \\
        &\Longrightarrow q = p \text{ by cancellation } \\
        &\Longrightarrow Q = P \\
        &\Longrightarrow n_p = 1. \ \square
    \end{align*}

    \item If $G$ is a group of order $45$, determine (with proof) $G$ up to isomorphism.\\
    \textbf{Proof:}Let $G$ be a group with $|G| = 45 = 3^2 \cdot 5$. Then, by Sylow's Theorem, we have
    \begin{align*}
        n_3 \equiv 1(\mod{3}) \text{ and } n_3 \big| 5 &\Longrightarrow n_3 = 1 \text{ and }\\
        n_5 \equiv 1(\mod{5}) \text{ and } n_5 \big| 3^2 = 9 &\Longrightarrow n_5 = 1
    \end{align*}
    Now, let $P \in Syl_3(G)$ and $Q \in Syl_5(G)$. Then $|P| = 3^2 = 9$ and $|Q|= 5$. Then, since $P$ and $Q$ are both finite subgroups of $G$, we have that 
    \[
        |PQ| = \dfrac{|P||Q|}{|P \cap Q|}.
    \]
    Now, note that $P \cap Q \leq P$ and $P \cap Q \leq Q$. Then, by La Grange's Theorem, we have
    \begin{align*}
        |P \cap Q| \big| |P| &\Longrightarrow |P \cap Q| \in \{1,3,9\} \text{ and } \\
        |P \cap Q| \big| |Q| &\Longrightarrow |P \cap Q| \in \{1,5\}
    \end{align*}
    Then we must have that $|P \cap Q| = 1$, so 
    \[
        |PQ| = \dfrac{|P||Q|}{|P \cap Q|} = \dfrac{9 \cdot 5}{1} = 45 = |G|.
    \]
    Now, since $n_3 = n_5 = 1$, we must have that $P,Q \trianglelefteq G$. Further, since $|PQ|=|G|$, we must have that $PQ = G$. Then, by the Normal Factoring Lemma, we must have that $G \cong P \times Q$.

    Now, note that $|P| = 9 \Longrightarrow P \cong \mathbb{Z}_9$ or $P \cong \mathbb{Z}_3 \times \mathbb{Z}_3$. Similarly, $|Q| = 5 \Longrightarrow Q \cong \mathbb{Z}_5$. Then we have 
    \[
        G \cong \mathbb{Z}_3 \times \mathbb{Z}_3 \times \mathbb{Z}_5 \text{ or } G \cong \mathbb{Z}_9 \times \mathbb{Z}_5. \ \square
    \]

    \item If $G$ is a group of order $105=3 \cdot 5 \cdot 7$, prove that $G$ has a cyclic subgroup of order $35$ and that this subgroup is normal in $G$.
    
    \item Let $p$ and $q$ be distinct primes. Prove that a group of order $p^2q$ is not simple.\\
    \textbf{Proof:} Let $G$ be a group such that $|G| = p^2q$ with $p$ and $q$ distinct primes. Then, by Sylow's Theorems, $\exists P \in Syl_p(G)$ and $\exists Q \in Syl_q(G)$. First, we investigate $Syl_p(G)$. 

    Note that, again by Sylow's Theorems, $n_p \equiv 1(\mod{p})$ and $n_p \big| q$. Then we must have that $n_p \in \{1,q\}$. Now, if $n_p = 1$, then $p \in Syl_p(G)$ is the sole sylow $p$-subgroup of $G$ and thusly, $P \trianglelefteq G$. Then $G$ is not simple and we are done. Then assume that $n_p = q$. Next, we investigate $Syl_q(G)$. 

    Again citing Sylow's Theorems, $n_q \equiv 1(\mod{q})$ and $n_q \big| p^2$. Then $n_q \in \{1,p,p^2\}$. Similarly to above, if $n_q = 1$, then $Q \in Syl_q(G)$ is the sole sylow $q$-subgroup and is thusly normal in $G$. Then we are done. So assume that $n_q \neq 1$. Then $n_q \in \{p,p^2\}$. Now, since $n_p = q$, we have
    \begin{align*}
        n_p = q &\Longrightarrow q \equiv 1(\mod{p}) \\
        &\Longrightarrow q \in \{1,p+1,2p+1, \dots \} \\
        &\Longrightarrow p < q.
    \end{align*}

    Assume $n_q = p$. Now, since $n_q \equiv 1(\mod{q})$, we must have that $n_q \in \{1, q+1, 2q+1, \dots \}$, so $p \in \{1,q+1,2q+1, \dots \}$. But $1 < p < q$, so we have a contradiction (since $p < q \Longrightarrow p < q+1$).

    Then assume that $n_q = p^2$. Now, since $|Q| = q$ and there are $p^2$ sylow $q$-subgroups in $G$, we must have that the total number of elements (not counting the identity) that are accounted for in the sylow $q$-subgroups is $p^2(q-1) = p^2q-p^2$. This means that there are exactly $p^2$ elements of $G$ remaining to be accounted for, all of which must be elements of $P \in Syl_p(G)$. Since $P \in Syl_p(G) \Longrightarrow |P| = p^2$, we must only have one sylow $p$-subgroup of $G$. That is, $n_p = 1$. Then $P \trianglelefteq G$ and $G$ is not simple. $\square$
    
    \item Prove that a group of order $132$ cannot be simple.\\
    \textbf{Proof:} Let $G$ be a group with order $132 = 2^2 \cdot 3 \cdot 11$. Then, by Sylow's Theorem, we have
    \begin{align*}
        n_2 &\equiv 1(\mod{2}) \text{ and } n_2 \big| 33 \\
        &\Longrightarrow n_2 \in \{1,3,5,7,9,11\}, \\
        n_3 &\equiv 1(\mod{3}) \text{ and } n_3 \big| 44 \\
        &\Longrightarrow n_3 \in \{1,4,22\}, \text{ and }\\
        n_{11} &\equiv 1(\mod{11}) \text{ and } n_{11} \big| 12 \\
        &\Longrightarrow n_{11} \in \{1,12\}.
    \end{align*}
    We will start with $n_{11}$. Now, if $n_{11}=1$, then $|Syl_{11}(G)| = 1$ and $P \in Syl_{11}(G) \Longrightarrow P \trianglelefteq G$ and we are done. Then assume $n_{11}=12$. Note that for each $Q_i \in Syl_{11}(G)$, $|Q_i| = 11$. Further, we have that, since the order of each $Q_i$ is prime and each $Q_i$ is unique, $Q_n \cap Q_m = \{e\} \ \forall n,m \in \{1,2, \dots, 12\}$. Then there are $10 \cdot 12 = 120$ unique elements from $G$ accounted for in $\bigcup_{i =1}^{12}{Q_i}$, $Q_i \in Syl_{11}(G)$, which leaves $12$ elements remaining.

    Now, we consider $n_3$. Since there are only $12$ distinct elements remaining in $G$, we know $n_3 \neq 22$. Further, if $n_3=1$ then $K \in Syl_3(G) \Longrightarrow K \trianglelefteq G$ and we are done. Then assume that $n_3 = 4$. Similarly to above, $Syl_3(G)$ will contain $4 \cdot 2 = 8$ unique elements between it's $4$ subgroups, leaving $4$ unique elements remaining in $G$.

    Lastly, we have $n_2$. Note that we must have that $n_2=1$ since $H \in Syl_2(G) \Longrightarrow |H| = 4$. Then $H \trianglelefteq G$ and $G$ is not simple. $\square$

    \item Let $P \in Syl_p(G)$ and let $N$ be a normal subgroup of $G$. Prove that $P \cap N$ is a Sylow $p$-subgroup of $N$.
    \item Let $G$ be a group of order $231 = 3 \cdot 7 \cdot 11$. Prove that $Z(G)$ contains a Sylow $11$-subgroup. Hint: One approach is to use the $\sfrac{N}{C}$ Theorem which states that if $H$ is a subgroup of $G$, then $\sfrac{N_G(H)}{C_G(H)}$ is isomorphic to a subgroup of $Aut(H)$.
    \item Prove that there is no simple group of order $300 = 2^2 \cdot 3 \cdot 5^2$. Hint: Assume that $G$ is not simple. By considering $n_5$, obtain a contradiction by showing $G$ is isomorphic to a subgroup of $S_6$.
\end{enumerate}

\newpage

\textbf{COMPREHENSIVE EXAM}\\
ALGEBRA 
Sprint 2026

\begin{enumerate}
    \item Let $G$ be a finite group and let $p$ be the smallest prime dividing the order of $G$. If $H$ is a subgroup of $G$ of index $p$, prove that $H$ is a normal subgroup of $G$.
    \item Let $G$ be a group of order $p^{\alpha}$ where $p$ is a prime and $\alpha$ is a positive integer.
    \begin{enumerate}
        \item Prove that $|Z(G)|>1$, where $Z(G)$ is the center of $G$.\\
        \textbf{Proof:} Let $G$ be a group of order $p^{\alpha}$ with $p$ a prime and $\alpha \in \mathbb{Z}^+$. We define $Z(G) = \{g \in G: gx = xg \ \forall x \in G\}$.  

        Now, by the class equation, 
        \[
            |G| = |Z(G)| + \sum[G:C_G(x)]
        \]
        where the sum is taken over all conjugacy classes of size greater than $1$ with only $1$ $x$ per conjugacy class. Note that we only take the sum over conjugacy classes of size greater than one because if the conjugacy class only has one element, then that element must commute with all elements of $G$ and is thusly in the center (which is already being counted by $|Z(G)|$ in the class equation).

        Now, by the Orbit-Stabilizer Theorem, $[G:C_G(x)]= \dfrac{|G|}{|C_G(x)|}$. Then we have
        \begin{align*}
            |G|=|Z(G)|+\sum[G:C_G(x)] &\Longrightarrow |G|=|Z(G)|+\sum{\dfrac{|G|}{|C_G(x)|}}\\
            &\Longrightarrow p^{\alpha} = |Z(G)|+\sum{\dfrac{p^{\alpha}}{|C_G(x)|}}
        \end{align*}
        Now, note that by La Grange's Theorem, $H \leq G \Longrightarrow |H| \big||G|$. Further, $C_G(x) \leq G$. Then we must have that $|C_G(x)| \big| |G|$, and thusly $|C_G(x)| = p^{l_x}$ with $1 \leq l_x < \alpha$. Then
        \begin{align*}
            p^{\alpha} = |Z(G)|+\sum{\dfrac{p^{\alpha}}{|C_G(x)|}} &\Longrightarrow p^{\alpha} = |Z(G)|+\sum{\dfrac{p^{\alpha}}{p^{l_x}}} \\
            &\Longrightarrow p^{\alpha} = |Z(G)|+\sum{p^{\alpha-l_x}}\\
            &\Longrightarrow p^{\alpha}-\sum{p^{\alpha-l_x}} =|Z(G)|\\
            &\Longrightarrow p(p^{\alpha-1}-\sum{p^{\alpha-l_x-1}}) = |Z(G)| \\
            &\Longrightarrow p \big| |Z(G)| \\
            &\Longrightarrow |Z(G)| > 1. \ \square
        \end{align*}
        \item Prove inductively that any group of order $p^{\alpha}$ contains a subgroup of order $p^{\alpha-1}$ for all $\alpha \geq 1$. You may use Cauchy's Theorem for abelian groups without proving it.\\
        \textbf{Proof:} Let $G$ be a group of order $p^{\alpha}$ with $p$ a prime and $\alpha \in \mathbb{Z}^+$. First, assume $\alpha = 1$. Then $|G| = p$ and we consider the trivial subgroup, $H = \{e\}$. Note that $|H| = 1 = p^{\alpha-1} = p^0$. Then our base case is true.

        Now, assume that $|G| = p^{k}, k < \alpha \Longrightarrow \exists H \leq G$ with $|H| = p^{k-1}$ is true. We need to show that $|G| = p^{k+1} \Longrightarrow \exists H \leq G$ with $|H| = p^{k}$ is true. Then assume that $|G| = p^{\alpha}$. Note that, by our proof in part $(a)$ above, $p \big| |Z(G)|$. Then, by Cauchy's Theorem for abelian groups, $\exists N \leq Z(G)$ such that $|N| = p$. Further, since $N \leq Z(G)$, $N \trianglelefteq G$ (any subgroup of the center of a group must be a normal subgroup of the entire group).

        Next, consider the quotient group $\sfrac{G}{N}$. Note that, since $G$ and $N$ are finite, 
        \[
            \big| \sfrac{G}{N} \big| = \sfrac{|G|}{|N|} = \dfrac{p^{\alpha}}{p} = p^{\alpha-1}.
        \]
        Then, by our inductive hypothesis, $\exists K \leq \sfrac{G}{N}$ such that $|K| = p^{\alpha-2}$. Then, by the Fourth Isomorphism Theorem, $\exists H \leq G$ such that $N \subseteq H$ and $\sfrac{H}{N} \cong K$. Then we have
        \begin{align*}
             |H| &= \big| \sfrac{H}{N} \big| \cdot |N| \\
             &= |K| \cdot |N| \\
             &= p^{\alpha-2} \cdot p \\
             &= p^{\alpha -1} \\
             &\Longrightarrow \exists H \leq G: |H| = p^{\alpha - 1}. \ \square
        \end{align*}
    \end{enumerate}
    \item Let $G$ be a group. A subgroup $H$ of $G$ is a \textit{characteristic subgroup} if $\phi(H)=H$ for all isomorphisms $\phi:G \longrightarrow G$.
    \begin{enumerate}
        \item If $H$ is a characteristic subgroup of $G$, prove that $H$ is a normal subgroup of $G$.
        \item Prove that $Z(G)$ is a characteristic subgroup of $G$.
    \end{enumerate}
    \item Let $G$ be a group of order $p^3qr$, where $p,q$, and $r$ are primes and $p>q>r$. Suppose that $P$ and $Q$ are distict Sylow $p$ subgroups with $|P \cap Q| > 1$.
    \begin{enumerate}
        \item Prove that $|P \cap Q| = p^2$.\\
        \textbf{Proof:} Let $P,Q \in Syl_p(G)$ be distinct. Let $|G| = p^3qr$ with $p,q,$ and $r$ primes, $p>q>r$, and $|P \cap Q| > 1$. 
        
        First, note that $|P|=|Q| = p^3$. Also note that $P \cap Q \leq P$ and $P \cap Q \leq Q$. Then, by La Grange's Theorem, $|P \cap Q| \big |P|$ and $|P \cap Q| \big |Q|$. Then $|P \cap Q| \in \{p,p^2,p^3\}$. Further, since $P$ and $Q$ are distinct, we know that $\exists a \in P$ such that $a \notin Q$. Similarly, $\exists b \in Q$ such that $b \notin P$. Then we must have that $P \cap Q$ is a proper subgroup of $P$ and also a proper subgroup of $Q$. Then we must have that $|P \cap Q| < p^3$, which means that $|P \cap Q| \in \{p,p^2\}$.
        
        Now, assume that $|P \cap Q| = p$. Then, since $P$ and $Q$ are both finite, we have
        \begin{align*}
            |PQ| &= \dfrac{|P||Q|}{|P \cap Q|} \\
            &= \dfrac{p^3 \cdot p^3}{p} \\
            &= p^5.
        \end{align*}
        Note that $PQ \leq G$, so $|PQ| \leq |G|$. Then we must have that $p^5 \leq p^3qr \Longrightarrow p^2 \leq qr$. Now, say $p = 5$, $q = 3$, and $r = 2$. Then $p^2 = 25$ and $qr = 6$. But this is a contradiction, which means that we would need for $p < q$ or $p < r$, which is a contradiction. Then we cannot have that $|P \cap Q| = p$. This means that the only remaining possibility is $|P \cap Q| = p^2$. $\square$
        \item Prove that $N_G(P \cap Q)$ contains both $P$ and $Q$. Here $N_G(P \cap Q)$ is the normalizer of $P \cap Q$ in $G$.
        \item Prove that $P \cap Q$ is normal in $G$.
    \end{enumerate}
    \item Let $G$ be a group of order $2 \cdot 3^2 \cdot 11 = 198$.
    \begin{enumerate}
        \item Prove that $G$ has a unique $11$-Sylow subgroup. \\
        \textbf{Proof:} Note that $G$ has order $11^1 \cdot 18$. Then, by Sylow's Theorem, $\exists P \in Syl_{11}(G)$, so $n_{11} \geq 1$. Further, Sylow's Theorem also gives us that $n_{11} \equiv 1(\mod{11})$ and $n_{11} \big| 18$. Then we must have that $n_{11}=1$, so $P$ is unique. $\square$
        \item Let $P$ be the unique $11$-Sylow subgroup of $G$. Prove $C_G(P)$ is a normal subgroup of $G$ where $C_G(P)$ is the centralizer of $P$ in $G$.\\
        \textbf{Proof:} Let $P \in Syl_{11}(G)$ with $n_{11}=1$, and let $C_G(P) = \{c \in G: cp=pc \ \forall p \in P\}$ be the centralizer of $P$ in $G$. Note that, since $n_{11}=1$, we must have that $P \trianglelefteq G$. Then we need to show that $gC_G(P)g^{-1} = C_G(P) \ \forall g \in G$. That is, we need to show that $gcg^{-1} \in C_G(P), \ c \in C_G(P) \ \forall g \in G$. Let $p \in P$. We have 
        \begin{align*}
            p(gcg^{-1}) &= gg^{-1}p(gcg^{-1}) \\
            &= g(g^{-1}p)gcg^{-1} \\
            &= g(p_1g^{-1}g)cg^{-1} \text{ since } P \trianglelefteq G\\
            &= gp_1cg^{-1} \\
            &= gcp_1g^{-1} \\
            &= (gcg^{-1})p \text{ since we established } p_1g^{-1} = g^{-1}p\\
            &\Longrightarrow gcg^{-1} \in C_G(P) \\
            &\Longrightarrow gC_G(P)g^{-1}=C_G(P) \\
            &\Longrightarrow C_G(P) \trianglelefteq G. \ \square
        \end{align*}
        \item Prove that $G$ has a normal $3$-Sylow subgroup.\\
        \textbf{Proof:} Let $G$ be defined as above. Then, since $|G| = 3^2 \cdot 22$, Sylow's Theorem tells us that $\exists Q \in Syl_3(G)$. Further, Sylow's Theorem also tells us that $n_3 \equiv 1(\mod{3})$ and $n_3 \big| 22$. Then we must have that $n_3 =1$. Then $Syl_3(G) = \{Q\}$.\\
        Next, note that, since $Syl_3(G) = \{Q\}$ and $Q \leq G$, we have that $gQg^{-1} \leq G$ $\forall g \in G$. Next, let $\phi:Q \Longrightarrow gQg^{-1}$ be defined by $\phi(q)=gqg^{-1}$. We aim to show that this map is a bijection. Let $\phi(q_1)=\phi(q_2)$. Then we have
        \begin{align*}
            \phi(q_1)=\phi(q_2) &\Longrightarrow gq_1g^{-1} = gq_2g^{-1} \\
            &\Longrightarrow gq_1g^{-1}g = gq_2g^{-1}g \\
            &\Longrightarrow gq_1=gq_2 \\
            &\Longrightarrow g^{-1}gq_1 = g^{-1}gq_2 \\
            &\Longrightarrow q_1=q_2,
        \end{align*}
        so $\phi$ is one-to-one. Next, let $x \in gQg^{-1}$. Then we have
        \begin{align*}
            x \in gQg^{-1} &\Longrightarrow gqg^{-1} = x \text{ for some } q \in Q \\
            &\Longrightarrow gqg^{-1}g = xg \\
            &\Longrightarrow gq = xg \\
            &\Longrightarrow g^{-1}gq = g^{-1}xg \\
            &\Longrightarrow q = g^{-1}xg,
        \end{align*}
        so we must have that $\phi(g^{-1}xg)=q$. Then $\phi$ is a bijection and we must have that $|gQg^{-1}| = |Q|$. Then we must have that $gQg^{-1} = Q$, and thusly $Q \trianglelefteq G$. $\square$
    \end{enumerate}
\end{enumerate}

\end{document}